\documentclass[11pt,a4paper]{article}
\usepackage[utf8]{inputenc}
\usepackage[T1]{fontenc}
\usepackage[ngerman]{babel}
\usepackage[margin=2.4cm]{geometry}
\usepackage{booktabs}
\usepackage{longtable}
\usepackage{xcolor}
\usepackage[colorlinks=true,urlcolor=blue!60!black,linkcolor=black]{hyperref}
\usepackage{parskip}
\usepackage{enumitem}
\setlist{nosep,leftmargin=*}

\newcommand{\code}[1]{\texttt{\small #1}}
\newcommand{\ok}{\textcolor{green!45!black}{\textbf{erledigt}}}
\newcommand{\offen}{\textcolor{orange!75!black}{\textbf{offen}}}
\newcommand{\upstream}{\textcolor{blue!60!black}{\textbf{upstream}}}

\title{AWI-ESM3-4-2-veg-HR, CMIP7 piControl\\[2mm]
\large Was seit der dritten Antwort passiert ist, gemessen an \code{cli117}}
\author{Jan Streffing, AWI}
\date{20. August 2026}

\begin{document}
% babel-ngerman macht " aktiv, was in Code-Beispielen zuschlaegt. Die Shorthands
% werden erst bei \begin{document} scharfgeschaltet, also muss das hierhin.
\shorthandoff{"}
\maketitle

\section{Kurzfassung}

Seit der dritten Antwort ist ein weiterer voller Lauf durch, \code{cli117}. Er
traegt neun Commits mehr als \code{cli116}, darunter zwei echte Fehler von uns und
das Gitter, auf das wir gewartet haben.

\begin{center}
\begin{tabular}{lrrr}
\toprule
& \textbf{\code{cli116}} & \textbf{\code{cli117}} & \textbf{Differenz} \\
\midrule
Dateien & 534 & 534 & $\pm 0$ \\
Reports & 534 & 534 & $\pm 0$ \\
HIGH    &  35 &  24 & $-11$ \\
MEDIUM  & 305 & 306 & $+1$ \\
\bottomrule
\end{tabular}
\end{center}

Die $-11$ verdecken, was tatsaechlich passiert ist. \textbf{27 HIGH sind
verschwunden}, und \textbf{16 neue sind dazugekommen, alle aus einer einzigen
Ursache}, die nicht bei uns liegt und sich von selbst erledigt. Beides steht unten.

\section{Zwei Fehler, die wirklich unsere waren}

Beide kamen ans Licht, weil ich deine Befunde nach Zustaendigkeit sortiert habe und
zwei davon in der falschen Schublade lagen.

\textbf{\code{slthick} lag auf dem falschen Gitter.} Die Regel las
\code{atm\_remapped\_1m\_lsm}, also die auf $640 \times 1296$ regriddete
Land-See-Maske, deklarierte aber \code{g122}, das reduzierte Gauss-Gitter mit
421120 Zellen. Von allen unseren Regeln mit einer \code{atm\_remapped\_*}-Eingabe war
sie die einzige mit \code{g122}, alle anderen tragen \code{g129}. Richtig ist
trotzdem \code{g122}: \code{slthick} liefert die Schichtdicken zu \code{mrsol} und
\code{tsl}, und wer die zusammen rechnet, braucht sie auf demselben Gitter. Die
Regel liest jetzt ein natives Feld.

Darunter lag ein zweiter Fehler, der den ersten verdeckt hat. Unser Schritt
broadcastete die vier Schichtdicken nur ueber die \emph{letzte} horizontale
Dimension. Auf einem unstrukturierten Gitter gibt es nur eine, dort ist das nie
aufgefallen. Auf dem regulaeren wurde daraus \code{lon}, und die Breitenachse
verschwand kommentarlos. \code{cli116} hat \code{slthick(depth, lon)} ausgeliefert.

\textbf{\code{depth} trug weiter \code{units = "meter"}} bei \code{tsl},
\code{mrsol} und \code{mrsll}. Der Grund ist derselbe Namensbruch wie damals bei
\code{plev}: unser Pass sucht die Achsen unter ihrem Namen aus der Datenanfrage, die
Datei fuehrt sie unter dem \code{out\_name}. Fuer \code{lev}, \code{plev} und
\code{height} sind die gleich, fuer \code{sdepth} nicht, das steht als \code{depth}
auf der Platte.

Eine Falle dabei, die der naheliegende Fix gestellt haette: \code{mrsol} gibt es
zweimal, die Schichtvariante mit \code{meter} und die \code{d10cm}-Variante mit
\code{cm}. Euer Checker erlaubt beide. Die CV-Einheit einfach durchzuschreiben
haette Zentimeter zu Metern erklaert und die Schicht um Faktor 100 verschoben. Wir
gleichen jetzt nur noch die Schreibweise an, nie die Einheit selbst.

\section{Nebenbei: unsere Koordinatentabelle war zwoelf Tage alt}

Beim Aufraeumen ist mir aufgefallen, dass wir gegen \code{CMIP7\_coordinate.json}
vom 2026-07-09 cmorisiert haben, waehrend \code{aicc} und \code{wcrp\_cmip7} gegen
den Stand vom 2026-07-21 validieren. Sieben Achsen weichen ab. Zwei davon sind
sichtbar geworden:

\begin{itemize}
\item \code{typetreebd/be/nd/ne} trugen alle vier \code{value = "trees"}. Das ist
  genau dein Befund aus der Runde zu \code{cli112}. Ich hatte das im Code
  ueberschrieben, weil ich annahm, die Tabelle sei eben so. Upstream hatte die
  spezifischen CF-Flaechentypen laengst nachgezogen, mein Workaround war eine
  Kruecke fuer eine veraltete Datei.
\item \code{deltasigt} hatte einen ganzen Absatz Prosa als \code{long\_name},
  inklusive literalem \verb|\n\n|. In jeder \code{mlotst}-Datei stand deshalb ein
  \code{long\_name} mit doppeltem Leerzeichen und ein \code{comment} mit dem Rest.
\end{itemize}

Die Tabelle ist jetzt auf dem Stand, gegen den auch geprueft wird, und ein Test
schlaegt an, wenn ein kuenftiger Refresh wieder eine aeltere einzieht.

\section{Das Zonalgitter, und ein Zeitfenster von drei Stunden}

Die vier Dateien ohne eigenes Gitter (\code{msftm} auf Tiefe und auf Dichte,
\code{hfbasin}, \code{sltbasin}) haben ihr Label bekommen. Registriert als
\code{g239}, gemergt am 18.08. Die Achsen passten schon vorher: 180 Baender von
$-89{,}5$ bis $89{,}5$ mit Schritt 1. Es fehlte nur das Label, vorher stand dort
\code{g130} und behauptete damit die 3146761 FESOM-Knoten.

Das hat 15 der 27 verschwundenen HIGH gebracht. Und alle 16 neuen, denn
\code{g239} kennt kein Checker:

\begin{center}
\begin{tabular}{rl}
\toprule
4 & \code{[ATTR004]} \code{grid\_label} Vokabularpruefung \\
4 & \code{[FILE001]} DRS-Verzeichnis \\
4 & \code{[FILE001]} DRS-Dateiname \\
4 & \code{[AICC002]} \code{grid\_label='g239' is not registered} \\
\bottomrule
\end{tabular}
\end{center}

Ich bin der Kette nachgegangen, weil mich interessiert hat, ob wir etwas falsch
gemacht haben:

\begin{center}
\begin{tabular}{lll}
\toprule
\textbf{Stufe} & \textbf{\code{g239}} & \textbf{Zeitpunkt} \\
\midrule
EMD \code{src-data}                & da      & 18.08., PR \#1294 \\
EMD \code{production}              & da      & \textbf{18.08. 08:22} \\
\code{CMIP7-CVs} \code{grid\_label} & fehlt   & letzter Sync \textbf{18.08. 05:34} \\
esgvoc \code{cmip7@1.2.18}          & fehlt   & veroeffentlicht 19.08. \\
\bottomrule
\end{tabular}
\end{center}

Der Sync von EMD in die CV laeuft als automatischer PR
(\code{.github/scripts/sync\_emd\_entries.py}). Der letzte gemergte war \#576 am
18.08. um 05:34. \code{g239} kam um 08:22 auf \code{production}, also zwei Stunden
und achtundvierzig Minuten zu spaet fuer genau diesen Lauf. \code{g238} hatte es
rechtzeitig geschafft und ist drin.

Der naechste Sync liegt schon als Entwurf vor, \code{CMIP7-CVs} \#592, und
\code{grid\_label/g239.json} ist darin enthalten. Es haengt also an einem Merge und
einer CV-Version danach. Nichts davon ist kaputt, und wir drehen \code{g239} nicht
zurueck: lieber 16 Befunde aus einem Terminproblem als vier Dateien, die ein Gitter
behaupten, auf dem sie nicht liegen.

Falls du an der Stelle mehr Einblick hast als ich: aus Sicht eines Datenproduzenten
ist das unangenehm, weil zwischen "Gitter ist registriert" und "Gitter validiert"
ein Fenster von Tagen liegt, in dem ein Produktionslauf HIGH-Befunde einsammelt,
die niemand beheben kann.

\section{Was uebrig bleibt}

Acht HIGH, keiner davon neu:

\begin{center}
\begin{tabular}{rl}
\toprule
5 & \code{plev7h}, unsere Modellausgabe schreibt die falschen Druckflaechen \\
1 & \code{pfull}, die Klimatologie war in diesem Lauf schon aktuell \\
1 & \code{vegtype}, der Tabellenwiderspruch \\
1 & \code{basin}, \code{cc-plugin-wcrp\#67} \\
\bottomrule
\end{tabular}
\end{center}

Zu \code{vegtype} noch ein Nachtrag, weil ich es diesmal nachgesehen habe: der
geforderte Term ist \code{needleleaf\_evergreen\_tree}, im Singular, waehrend seine
drei Geschwister im Plural stehen und wir \code{needleleaf\_evergreen\_trees}
schreiben. Der Singular steht \emph{in beiden} Koordinatentabellen, auch in der
offiziellen vom 2026-07-21. Es ist also kein Versaeumnis von uns und auch keine
Altlast unserer Kopie.

Der eine MEDIUM mehr ist ein CF-2.4-Befund auf \code{slthick}, und den handelt man
sich genau dadurch ein, dass die Datei jetzt eine echte horizontale Dimension hat.
Gleiche Klasse wie die anderen 75.

\section{Die fuenf fehlenden Sidecars aus \code{cli116}}

Beim Nachzaehlen von \code{cli116} lagen nur 529 Reports auf 534 Dateien. Fuenf
Dateien hatten gar keinen Sidecar. Die Ursache lag nicht an den Dateien.

Die fuenf sind vollstaendig und lassen sich einwandfrei
pruefen, ich habe genau denselben Aufruf spaeter von Hand wiederholt und vier davon
kamen ohne einen einzigen Befund zurueck.

Es lag an \code{cchecker.py} selbst. Er stirbt sporadisch schon beim Import, bevor er
ueberhaupt eine Datei anfasst:

\begin{quote}
\code{FileNotFoundError: [Errno 2] No such file or directory}\\
\code{~~File ".../bin/cchecker.py", line 10, in <module>}\\
\code{~~~~from compliance\_checker.cf.util import download\_cf\_standard\_name\_table}\\
\code{~~File ".../compliance\_checker/cf/util.py", line 17, in <module>}
\end{quote}

Unser Verdacht: parallele Shards, die gleichzeitig den Cache der
CF-Standardnamentabelle anlegen wollen, und einer sieht die Datei halb geschrieben.
Dazu passt, dass es in demselben Lauf und derselben Umgebung 529 mal geklappt hat
und fuenf mal nicht, und dass der Wiederholungsversuch von Hand jedes Mal
durchlief. Reproduzierbar ist es fuer mich nicht, deshalb bleibt es ein Verdacht.
Falls du das kennst, wuerde mich interessieren woran es bei euch lag.

\section{Warum das lautlos war, und was jetzt dagegen laeuft}

Das ist derselbe lautlose Ausfall, den ich in der dritten Antwort im Zusammenhang mit
\code{CMIP7\_TABLES\_PATH} beschrieben habe, nur mit einer anderen Ursache: der
QC-Schritt hat protokolliert, dass kein JSON entstand, und ist weitergelaufen. Eine
Datei ohne Sidecar wird danach nirgends mehr gezaehlt. Der Lauf sieht dadurch
sauberer aus als er ist, und das ist die unangenehme Richtung des Fehlers.

Der QC-Schritt wiederholt den Aufruf jetzt bis zu dreimal, mit gestaffelter Pause,
damit zwei kollidierte Shards nicht gleich wieder kollidieren. Wiederholt wird nur,
wenn ueberhaupt kein Report auf der Platte liegt. Ein Report mit Befunden und
Exitcode $\neq$ 0 ist der Normalfall und wird wie bisher durchgereicht. Bleibt es
auch nach dem letzten Versuch dabei, sagt die Fehlermeldung jetzt ausdruecklich,
dass diese Regel ohne Report dasteht.

\section{Was das fuer den Bestand heisst}

Fuer \code{cli116} sind es jetzt 534 Reports auf 534 Dateien, und \code{cli117}
hat von sich aus 534 auf 534 geliefert. Der Wiederholungsversuch hat dort nicht
einmal ausgeloest, es gab schlicht keinen Importfehler. Die Zahl ist also echt und
nicht gerettet. Die frueheren Laeufe haben wir nicht nachgezaehlt. Falls dir in Zahlen aus
\code{cli108} bis \code{cli115} etwas nicht aufgeht, koennte das derselbe Effekt
sein, und ich schaue gezielt nach.

\end{document}
